%; whizzy chapter -dvi
% -initex iniptex -latex platex -format platex -bibtex jbibtex -fmt fmt
% 以上 whizzytex を使用する場合の設定。
 
%     Tokyo Debian Meeting resources
%     Copyright (C) 2012 Junichi Uekawa
%     Copyright (C) 2011, 2015 Nobuhiro Iwamatsu

%     This program is free software; you can redistribute it and/or modify
%     it under the terms of the GNU General Public License as published by
%     the Free Software Foundation; either version 2 of the License, or
%     (at your option) any later version.

%     This program is distributed in the hope that it will be useful,
%     but WITHOUT ANY WARRANTY; without even the implied warranty of
%     MERCHANTABILITY or FITNESS FOR A PARTICULAR PURPOSE.  See the
%     GNU General Public License for more details.

%     You should have received a copy of the GNU General Public License
%     along with this program; if not, write to the Free Software
%     Foundation, Inc., 51 Franklin St, Fifth Floor, Boston, MA  02110-1301 USA

%  preview (shell-command (concat "evince " (replace-regexp-in-string "tex$" "pdf"(buffer-file-name)) "&"))

%%ここからヘッダ開始。

\documentclass[mingoth,a4paper]{jsarticle}
\usepackage{monthlyreport}
% 日付を定義する、毎月変わります。
\newcommand{\debmtgyear}{2015}
\newcommand{\debmtgmonth}{5}
\newcommand{\debmtgdate}{23}
% started from zero:
% (let ((year 2013) (month 7)) (+ (* (- year 2005) 12) month -1))
\newcommand{\debmtgnumber}{126}

\begin{document}

\begin{titlepage}
\thispagestyle{empty}
% タイトルページ:編集必要な部分は最初のマクロに飛ばすこと

\vspace*{-2cm}
第\debmtgnumber{}回 東京エリア Debian 勉強会資料\\
\hspace*{-2cm}
\includegraphics{image2012-natsu/dotdeb.pdf}\\
\hfill{}\debmtgyear{}年\debmtgmonth{}月\debmtgdate{}日

% ここはアップデートすること
% 全角文字にしないとフォントのサイズが合わないので注意
\rotatebox{10}{\fontsize{30}{30} {\gt 特集：自然言語処理チームとパッケージ}}\\

\vspace*{-2cm}
\hfill{}\includegraphics[height=6cm]{image200502/openlogo-nd.eps}
\end{titlepage}

\newpage

\begin{minipage}[b]{0.2\hsize}
 \definecolor{titleback}{gray}{0.9}
 \colorbox{titleback}{\rotatebox{90}{\fontsize{80}{80} {\gt デビアン勉強会} }}
\end{minipage}
\begin{minipage}[b]{0.8\hsize}
\hrule
\vspace{2mm}
\hrule
\begin{multicols}{2}
\tableofcontents
\end{multicols}
\vspace{2mm}
\hrule
\end{minipage}

\dancersection{事前課題}{野島 貴英}

今回の事前課題は以下です:
\begin{enumerate}
\item 本日、何の作業をやるかを宣言ください。
\item (オプション) どこで今回の勉強会の開催を知りましたか？
\item (オプション) 何について聞きたい／参加者と話をしたいですか？
\end{enumerate}
この課題に対して提出いただいた内容は以下です。
\begin{multicols}{2}
{\small
\begin{prework}{ $BLnEg(B }
  \begin{enumerate}
  \item Q.hack time$B$K2?$r$7$^$9$+!)(B\\
    A. Nook HD+$B$r$=$m$=$m(Bdebian$B$K!#(B
  \item ($B%*%W%7%g%s(B)Q.$B2?$K$D$$$FJ9$-$?$$!?;22C<T$HOC$r$7$?$$$G$9$+!)(B\\
    A. $B%=%U%H%&%'%"<+M3$r<i$jDL$9;v$KBP$7$F!"2?$,=PMh$k$+$K$D$$$F!#(B
  \end{enumerate}
\end{prework}

\begin{prework}{ wskoka }
  \begin{enumerate}
  \item Q.hack time$B$K2?$r$7$^$9$+!)(B\\
    A. MIPS Debian $B$N0\?"(B
  \item ($B%*%W%7%g%s(B)Q.$B$I$3$G:#2s$NJY6/2q$N3+:E$rCN$j$^$7$?$+!)(B\\
    A. $B$=$NB>(B
  \end{enumerate}
\end{prework}

\begin{prework}{ koedoyoshida  }
  \begin{enumerate}
  \item Q.hack time$B$K2?$r$7$^$9$+!)(B\\
    A. DDTSS,PyconJP$B4XO"(B
  \item ($B%*%W%7%g%s(B)Q.$B$I$3$G:#2s$NJY6/2q$N3+:E$rCN$j$^$7$?$+!)(B\\
    A. $BM'C#$dCN$j9g$$$+$iD>@\(B
  \end{enumerate}
  DDTSS: http://ddtp.debian.net/ddtss/index.cgi/ja
\end{prework}

\begin{prework}{ dictoss }
  \begin{enumerate}
  \item Q.hack time$B$K2?$r$7$^$9$+!)(B\\
    A. kFreeBSD$B$N(BIPsec$B4XO"$rD4$Y$k(B
  \item ($B%*%W%7%g%s(B)Q.$B$I$3$G:#2s$NJY6/2q$N3+:E$rCN$j$^$7$?$+!)(B\\
    A. Debian JP$B$N%a!<%j%s%0%j%9%H(B
  \end{enumerate}
\end{prework}

\begin{prework}{ yy\_y\_ja\_jp }
  \begin{enumerate}
  \item Q.hack time$B$K2?$r$7$^$9$+!)(B\\
    A. DDTSS 
  \item ($B%*%W%7%g%s(B)Q.$B$I$3$G:#2s$NJY6/2q$N3+:E$rCN$j$^$7$?$+!)(B\\
    A. $B$=$NB>(B
  \item ($B%*%W%7%g%s(B)Q.$B2?$K$D$$$FJ9$-$?$$!?;22C<T$HOC$r$7$?$$$G$9$+!)(B\\
    A. DDTSS
  \end{enumerate}
\end{prework}

}
\end{multicols}

\dancersection{Debian Trivia Quiz}{野島 貴英}

 Debianの昨今の話題についてのQuizです。

今回の出題範囲は\url{debian-devel-announce@lists.debian.org} や \url{debian-news@lists.debian.org}に投稿された
内容などからです。

\begin{multicols}{2}
%; whizzy-master ../debianmeetingresume201311.tex
% 以上の設定をしているため、このファイルで M-x whizzytex すると、whizzytexが利用できます。
%

\santaku
{Jessieが無事リリースされました。いつだったでしょうか？}
{2015/4/11}
{2015/4/18}
{2015/4/25}
{C}
{無事4/25にJessie(Debian 8)がリリースされました。Linux kernelは3.16.7、GNU Compiler Collection 4.9.2、GNOME 3.14、Apache 2.4.10、PHP 5.6.7などのバージョンのものが搭載されました。搭載されているパッケージ数は全部で43,000パッケージ以上もあります。}

\santaku
{Jessieで初めて追加されたものではないものが混ざっています。どれ？}
{Debian Games Blend}
{OpenJDK}
{androidsdk-tools}
{B}
{Jessieでの変更点はリリースノート(https://www.debian.org/releases/ jessie/amd64/release-notes/index.ja.html)にあります。Debian Gamesチームよりゲームの詰め合わせのDebian Blendが追加されたり、androidの開発ツールの一部も追加されたりしています。}

\santaku
{Debian GNU/Hurd 2015も2015/4/30にリリースされました。心臓部のGNU Machのバージョンはいくつ？}
{1.6}
{1.5}
{1.4}
{B}
{2015/4/10〜2015/4/15周辺でGNU/Hurd（upstream側)のアップデートが行われ、GNU Hurdはversion 0.6に、GNU Machは1.5になりました。未だにKVM/QEMUなどの仮想環境での動作が推奨の32bitシステムですが、多くのOSがクラウド環境で動作している世の中ですので、ホビー用途の利用でそろそろ注目されても良いかも？と思う次第です。}

\santaku
{2015年のGSoCに採択された、Debian MIPS portsについての開発内容は次のどれ？}
{多数のビルド出来ないパッケージを、ちゃんとビルドできるようにする。}
{新しいMIPS CPUへの対応}
{新しいMIPS CPU搭載製品への対応}
{A}
{Debian MIPS portsの多数のパッケージは再ビルドが出来ない、インストールに失敗するものが多数ある状況です。こちらを再ビルドできるように修正し、インストール出来るようにする事が採択されました。GSoCに限らず、いつでも本件の協力者募集中ですので、MIPS portsに興味ある方はパッケージの修正に是非ご協力くださいませー。}

\santaku
{遂にhttp.debian.netがdebian.orgのインフラに移動となりました。新しいURLはどれ？}
{http://http.debian.org/debian}
{http://httpredir.debian.org/debian}
{http://www.debian.org/}
{B}
{http.debian.netはHTTPのRedirectを活用し、ユーザに最も適切な場所にある、パッケージのミラー先を教えてくれるサービスです。遂に、実験的サービスの扱いであるdebian.netからDebian公式サービスであるdebian.orgへ移動となりました。http.debian.netをsource.listに設定している人は、httpredir.debian.orgへ修正をおねがいします。また、不具合は"mirrors" pseudo-packageでBTSするようです。参照：http://bugs.debian.org/mirrors　}

\santaku
{debianへのlibdvdcss/ZFSパッケージ搭載の件について、2015/5のBit From DPLで報告された状況は以下のどれ?}
{DPLがいろいろ議論して回っている状況}
{搭載日確定した}
{搭載を諦めた}
{A}
{2015/5/19現在、FTPMasters, SFLC, FSFと議論中とのこと。さらに近々Linux Foundationとも議論予定。現在、当初想定したよりも、実現に関する状況が複雑とのことで、本件いつ決着するか見えなくなったとの事です。}


\end{multicols}

\dancersection{最近のDebian関連のミーティング報告}{野島 貴英}

\subsection{第125回東京エリアDebian勉強会}

\begin{itemize}
\item 場所はスクウェア・エニックスさんのセミナルームをお借りしての開催でした。
\item 参加者は13名でした。
\item セミナ内容はまえだこうへいさんによる「.debからPythonパッケージへの変遷」でした。
\item 残りの時間でhack timeを行い、成果発表をしました。
\item 宴会の代わりに、「祥龍房 新宿イーストサイドスクエア店」で夕食会をやりました。
\end{itemize} 

 セミナはまえださんにより、PythonパッケージのローカルDebian CIに関する発表と、Pythonに関するDebian公式パッケージの考察が行われました。大統一Debian勉強会から、毎年発表のはや３年越しのシリーズものとなります。昨今の言語関係の流行りや事情から、言語のパッケージをDebianのパッケージ化を検討した方が良い場合と、あまり向かない場合がある状況とのことです。会場にHaskelの詳しい人、Perlに詳しい人がおり、様々に議論がなされました。
 
% % (query-replace-regexp "<.*?>" "")
% % (query-replace-regexp "^[	 ]\+" "")

%-------------------------------------------------------------------------------
\dancersection{自然言語処理チーム(pkg-nlp-ja)とパッケージ}{野首 貴嗣}
%-------------------------------------------------------------------------------

\subsection{NOTICE}
 私は専門家ではありません。

\subsection{pkg-nlp-ja}
 Debianにおける自然言語処理パッケージの管理
 \url{https://alioth.debian.org/projects/pkg-nlp-ja/}

\subsection{管理パッケージ}
\begin{itemize}
\item ChaSen
\begin{itemize}
  \item dirts (Double Array実装)
\end{itemize}
\item 辞書
\begin{itemize}
\item ipadic
\item naist-jdic
\item unidic
\end{itemize}
\end{itemize}

\subsection{その他のNLPパッケージ}
\begin{itemize}
\item 解析
  \begin{itemize}
  \item MeCab
  \item Juman
  \end{itemize}
\item 変換
  \begin{itemize}
  \item Anthy (pkg-anthy)
  \item KAKASI
  \item mozc
    \begin{itemize}
    \item 変換についてはpkg-imeで管理
    \end{itemize}
  \end{itemize}
\end{itemize}

\subsection{何ができるか}
形態素解析
\begin{itemize}
\item 形態素(単語)を調べる
\item 単語の属性情報を得る
  \begin{itemize}
  \item 品詞・読み等
    \begin{itemize}
    \item 辞書に依存
    \end{itemize}
  \end{itemize}    
\end{itemize}

 ChaSen, MeCab, Juman
 ref:\url{http://www.phontron.com/nlptools.php?lang=ja}

\subsection{応用例}
\begin{itemize}
\item 検索エンジン
  \begin{itemize}
  \item 転地インデックスの作成
  \item Groonga, Namazu等
  \end{itemize}
\item 音声合成 (Text-to-Speech)
  \begin{itemize}
  \item open-jtalk
  \end{itemize}n
\item 分かち書き
  \begin{itemize}
  \item word2vecの下処理
  \end{itemize}
\item かな漢字変換
\end{itemize}

\subsection{動作原理}

\begin{itemize}
\item 解析対象を確からしい単位で分割
  \begin{itemize}
  \item 例: 「東京都府中市」
    \begin{itemize}
    \item x 「東」「京都」「府」「中」「市」
    \item o 「東京」「都」「府中」「市」
    \end{itemize}
  \item 単純なルールでは処理できない
\end{itemize}

\subsection{解析の手がかり}

\begin{itemize}
\item マッチする単語の長さ
  \begin{itemize}
  \item 最長一致(KAKASI)
  \end{itemize}
  ref: KAKASIの実装と課題
\item 品詞情報の利用
  \begin{itemize}
    \item 前後に来やすい品詞のつながり
    \item  単語のつながり
    \item  単語の頻度/スコア
  \end{itemize}
\item 確率
  \begin{itemize}
  \item 出現率、連接確率
  \end{itemize}
\end{itemize}

\subsection{アルゴリズム}

\begin{itemize}
\item コスト最小法
  \begin{itemize}
  \item ChaSen
  \end{itemize}
\item CRF (条件付き確率場)
  \begin{itemize}
  \item MeCab
  \end{itemize}
\end{itemize}
ref: 日本語形態素解析入門(pdf)
\end{itemize}

\subsection{辞書探索}

 Common Prefix Search
\begin{itemize}
 \item 同じ接頭語を持つデータを高速に検索
\end{itemize}  

アルゴリズム
\begin{itemize}
\item トライ
 \begin{itemize}
  \item パトリシアトライ
 \end{itemize}
 \item Double Array
   \begin{itemize}
   \item dartsパッケージ
   \end{itemize}
\end{itemize}

\subsection{トライ(Trie)}

文字単位の木構造

\begin{verbatim}
 漢-+->字-+->化
   |　　　|
   |     +->語
   +->音
\end{verbatim}

ref: \url{http://ja.wikipedia.org/wiki/基数木}

\subsection{Double Array}

 トライ構造を2つの配列で表現

\begin{itemize}
\item BASE配列 (子ノード番号へのオフセット)
\item CHECK配列 (親ノード番号)
  \begin{itemize}
  \item 単純な状態遷移表だと疎な配列になる
  \item 構築に計算が必要
    \begin{itemize}
    \item 配列の空き要素を見つける必要がある
    \item 動的な更新には向かない
    \end{itemize}
  \end{itemize}
\end{itemize}
ref: ダブル配列の実装方法

\subsection{辞書}

\begin{itemize}
\item KAKASIDIC(kanwadic)
  \begin{itemize}
  \item SKKDICベース
  \end{itemize}
\item ipadic
  \begin{itemize}
  \item ChaSen, MeCab
  \end{itemize}   
\item naist-jdic
  \begin{itemize}
  \item ChaSen, MeCab
  \end{itemize}    
\item jumandic
  \begin{itemize}
  \item MeCab, JUMAN
  \end{itemize}    
\end{itemize}

\subsection{単語の追加: KAKASI}

辞書ファイルの作成(EUC-JP)
\begin{commandline}
# よみ [空白] 漢字
けいさんしょう 経産省
けいざいさんぎょうしょう 経済産業省
\end{commandline}
kakasiコマンドの引数に辞書のパスを追加
\begin{commandline}
$ echo '日本の経産省 ' | kakasi -w -iutf8 -outf8 ./extdic
日本 の 経産省
$ echo '日本の経産省 ' | kakasi -w -iutf8 -outf8
日本 の 経 産 省
\end{commandline}

\subsection{単語の追加: ChaSen}

\begin{itemize}
\item ipadicのソースを取得
\item 追加したい単語の意味に近いものを探す
\item そのエントリーをコピーし、スコアはそのままで単語を修正する
\item ./configure \&\& make 
\end{itemize}
 注: スコアを計算するツールが付属していない

\subsection{単語の追加: MeCab}

\begin{enumerate}
\item ChaSenと同じ方法
  \begin{itemize}
  \item 活用のある品詞はすべて開く必要がある
  \end{itemize}
\item mecab-cost-trainで正確なコストを算出
  \begin{itemize}
  \item 元となるコーパスを用意する必要がある
  \item cf: mecab-ipadic-neologd
  \item MeCab の IPA 辞書を再学習させてみる
  \item MeCabの辞書をカスタマイズする
  \end{itemize}
\end{enumerate}

\subsection{参考}

 自然言語処理ツール\\
\url{http://www.phontron.com/nlptools.php?lang=ja}

 日本語で読める自然言語処理のチュートリアルスライドまとめ\\
\url{http://blog.unnono.net/2015/04/nlp-tutorial.html}

\newpage

\vspace*{15cm}
\hrule
\vspace{2mm}
\includegraphics[width=2cm]{image200502/openlogo-nd.eps}
\noindent \Large \bf Debian 勉強会資料\\
\noindent \normalfont \debmtgyear{}年\debmtgmonth{}月\debmtgdate{}日 \hspace{5mm}  初版第1刷発行\\
\noindent \normalfont 東京エリア Debian 勉強会 （編集・印刷・発行）\\
\hrule

\end{document}
